\documentclass[11pt]{article}
\usepackage[margin=1in]{geometry}
\usepackage{booktabs}
\usepackage{longtable}
\usepackage{amsmath}
\usepackage{hyperref}
\hypersetup{colorlinks=true,linkcolor=blue,citecolor=blue,urlcolor=blue}

\title{The Political Economy of Local Land-Use Regulation\\
\large The theory line: fragmented regulation as a fiscal-federalism breakdown}
\author{Dan Post}
\date{August 2026}

\begin{document}
\maketitle

\section{Purpose}

The project studies the \emph{market for housing regulation}. Its empirical spine is an item-level
dataset of discretionary land-use decisions, built by parsing SF Planning Commission minutes and
being extended to the roughly 109 land-use regulators of the nine-county Bay Area, paired with the
by-right zoning envelope and with dated state-preemption exposure.

This directory is the \textbf{theory line}, not a data probe. It holds the model that tells the
empirical work what to measure and why: it defines the two regulatory margins the dataset is built
to separate, states the propositions the estimation is meant to test, and names the counterfactuals
the measurement is ultimately for. Its claim on the project is that a purely descriptive account of
land-use stringency cannot answer either question the project cares about --- how large the
fragmentation distortion is, and how much of state preemption localities claw back --- because both
are counterfactual, and both require an equilibrium object rather than a level.

\section{What is in this directory}

\begin{table}[h]
\centering
\begin{tabular}{@{}p{0.26\textwidth}p{0.32\textwidth}p{0.34\textwidth}@{}}
\toprule
File & What it is & How it was produced \\
\midrule
\texttt{toy\_model.tex} &
The model note, ``Fragmented Housing Regulation as a Fiscal-Federalism Breakdown'' (June 2026):
environment, two margins, homevoter objective, three propositions, identification, data mapping. &
Written by hand; the analytical core of the project. Compiles against
\texttt{toy\_model.bib} with \texttt{natbib}/\texttt{plainnat}. \\
\addlinespace
\texttt{toy\_model.pdf} & Compiled note. & \texttt{pdflatex} + \texttt{bibtex}. \\
\addlinespace
\texttt{toy\_model.bib} & Bibliography (Oates, Besley--Coate, Fischel, Glaeser, Hsieh--Moretti,
Bartik et al., Brueckner). & Hand-assembled. \\
\addlinespace
\texttt{guren\_meeting\_slides.tex} &
A 10-frame Beamer deck presenting the same framework for the Adam Guren meeting: the two margins,
the three-layer structure (demand / supply / regulation), the reaction function, and four explicit
questions for the advisor. & Written 2026-06-16 alongside the demand-data report. \\
\addlinespace
\texttt{guren\_meeting\_slides.pdf} & Compiled deck. & \texttt{pdflatex}. \\
\addlinespace
\texttt{.nav}, \texttt{.snm} & Beamer navigation auxiliaries. & Compilation by-products. \\
\bottomrule
\end{tabular}
\end{table}

\section{The model}

\subsection{Environment}

One metro (the Bay Area) contains $J$ jurisdictions $j = 1,\dots,J$, grouped into $S \ll J$ housing
submarkets $s(j)$, all drawing on a \emph{single} regional labor market. A worker employed anywhere
in the metro earns the regional wage $w_m = w(N)$ with $w' > 0$ and $N \equiv \sum_j N_j$, which
carries the agglomeration rent. This wage coupling is what makes the model's central asymmetry
structural rather than assumed: a household's income depends on the \emph{metro}, not on the
jurisdiction it sleeps in.

A household of type $\theta$ living in $j$ has preferences
\[
  U \;=\; c \;+\; \alpha \ln h \;-\; \gamma\, g(D_j), \qquad g' > 0,
  \qquad y(\theta) = w_m \theta ,
\]
where $D_j = N_j / \bar L_j$ is local density and $-\gamma g(D_j)$ is the \emph{local} disamenity of
density. Because low density enters as a normal good, high-$\theta$ residents are the natural
demanders of restriction. Within submarket $s$, market clearing gives the price index
\begin{equation}
  p_s \;=\; \frac{\alpha N_s}{H_s},
  \qquad N_s = \!\!\sum_{j:\, s(j)=s}\!\! N_j,
  \qquad H_s = \!\!\sum_{j:\, s(j)=s}\!\! H_j .
  \label{eq:price}
\end{equation}

\noindent Equation \eqref{eq:price} \emph{is} the externality. Housing built in any jurisdiction of
$s$ lowers the price faced by incumbents in \emph{every} jurisdiction of $s$, while density is local.
The cost of approval is borne locally; the benefit --- lower submarket prices, more regional
employment and hence a higher $w_m$ --- is shared. That is exactly the spillover condition the
Decentralization Theorem excludes, so decentralized provision of regulation is inefficient.

\subsection{Two margins of regulation}

The model's key structural move comes from an empirical fact visible in the minutes: among projects
that reach a vote, approval is nearly universal. Restriction does not operate through denial at the
hearing. It operates before and around it, on two margins.

\begin{itemize}
\item \textbf{The by-right envelope $\bar z_j$.} The zoning map's use and height/bulk districts
define what a project may obtain \emph{ministerially} --- with certainty, no hearing. Projects
inside the envelope never appear on the commission calendar. This is the margin that LLM-based
zoning measurement captures, and the margin state preemption targets.
\item \textbf{The discretionary friction $\tau_j$.} A project exceeding the envelope enters
discretionary review and is built only after paying a \emph{discretionary tax} whose components the
minutes record: process exposure (the share of activity forced off the by-right path), mobilized
opposition (the operative form of $g(D_j)$, measured by stance-marked speaker counts), conditioning
(approval-with-conditions), delay (continuances), and rare denial. This is the margin LLM zoning
measurement \emph{cannot} see and preemption leaves untouched.
\end{itemize}

\noindent Effective stringency is then $r_j = g(\bar z_j, \tau_j)$, increasing in the tightness of
the envelope (a smaller $\bar z_j$ pushes more activity into discretion) and in the discretionary
tax. The two margins are \emph{substitutes} in producing restriction but face different costs --- and
preemption acts on one and not the other. This definition is what makes leakage a coherent object;
it is the piece the descriptive literature lacks.

\subsection{The homevoter's problem}

Regulation is chosen by a vote-maximizing local authority answering to incumbent homeowners (share
$\lambda_j$, political weight $\phi_j$). The two margins map to two bodies --- the council sets the
slow map $\bar z_j$, the planning commission sets the fast discretionary tax $\tau_j$ --- which
resolves the usual unitary-authority shortcut. Taking neighbors' policies as given, an incumbent's
payoff is
\begin{equation}
  W^{\text{own}}_j
  \;=\; \underbrace{p_{s(j)}\big(\bar z_{-j}, \tau_{-j};\, \bar z_j, \tau_j\big)}_{\text{submarket asset value}}
  \;-\; \underbrace{\gamma\, g\!\big(D_j(\bar z_j, \tau_j)\big)}_{\text{local disamenity}},
  \label{eq:owner}
\end{equation}
where $-j$ is the profile of all other jurisdictions and only co-submarket neighbors ($s(k) = s(j)$)
enter with nonzero weight. The authority maximizes
\begin{equation}
  \Omega_j \;=\; \phi_j \lambda_j\, W^{\text{own}}_j \;+\; \big(1 - \phi_j \lambda_j\big) W^{\text{ent}}_j
  \qquad \text{over } (\bar z_j, \tau_j),
  \label{eq:obj}
\end{equation}
so $\phi_j \lambda_j$ is the model's homevoter weight and $W^{\text{ent}}$ the payoff of entrants.

\subsection{Solution concept}

Because $p_{s(j)}$ in \eqref{eq:owner} depends on \emph{neighbors'} policies, \eqref{eq:obj} is not a
decision problem but a game. The solution is the Nash equilibrium in policies, characterized by a
reaction function on each margin,
\[
  \big(\bar z_j, \tau_j\big) \;=\; R_j\big(\bar z_{-j}, \tau_{-j}\big),
  \qquad -j = \text{co-submarket neighbors only},
\]
and it is compared against the allocation chosen by an authority that internalizes the submarket (or
region) --- the planner who values entrants' access to $w_m$ and the cross-jurisdiction price
externality in \eqref{eq:price}.

\section{Preliminary results}

\subsection{The free-riding wedge}

The mechanism is an attenuation result. The effect of $j$'s own approvals on the price its
homeowners care about is diluted by submarket size:
\begin{equation}
  \frac{\partial p_{s(j)}}{\partial\, \text{supply}_j} \;=\; -\,\frac{p_s}{H_s}.
  \label{eq:attenuation}
\end{equation}
A jurisdiction that is small within its submarket bears the \emph{full} local disamenity of any
approval but captures only a $1/H_s$ share of the price relief; the rest leaks to co-submarket
neighbors. It therefore under-supplies and free-rides on their construction. Equation
\eqref{eq:attenuation} is the microfoundation of the fragmentation externality and drives everything
below.

\paragraph{Proposition 1 (Fragmentation over-restricts; the distortion is the Oates wedge).}
Holding $(\phi_j\lambda_j, \gamma, \bar L_j)$ fixed, a jurisdiction's privately optimal restriction
is \emph{increasing in the number of jurisdictions sharing its submarket}. As $j$ becomes atomistic
the price-relief motive vanishes, the local-disamenity motive dominates, $\bar z_j$ tightens and
$\tau_j$ rises. The \emph{distortion} is the gap between decentralized Nash regulation
$\{(\bar z_j^\ast, \tau_j^\ast)\}$ and the regulation an internalizing authority would choose. This
is the model's answer to Question 1, and it is precisely the inefficiency the Decentralization
Theorem's spillover condition rules in.

\paragraph{Proposition 2 (Regulation is strategically substitutable within a submarket).}
For co-submarket jurisdictions, if a neighbor loosens (raises $\bar z_k$ or lowers $\tau_k$), the
submarket price falls, which \emph{raises} a homevoter-pivotal $j$'s incentive to restrict: $j$
tightens. Neighbors' regulations are strategic substitutes, each free-riding on the others' supply,
and the Nash equilibrium delivers less regional construction than the internalized optimum.

\medskip
\noindent \textbf{Why Proposition 2 is the centerpiece.} The reaction-function slope is not merely an
empirical regularity; it is the direct test of the federalism breakdown. Under the Decentralization
Theorem's efficient case --- no spillovers --- jurisdictions would not respond to one another's
regulation at all, and the cross-jurisdiction slope would be zero. A robustly \emph{nonzero} slope is
the spillover made visible: its sign and magnitude establish that the Oates condition fails and
quantify how badly. This is also why the project must scale beyond San Francisco: single-jurisdiction
data cannot, even in principle, reveal a reaction function.

\paragraph{Proposition 3 (Leakage: substitution from envelope to discretion).}
State preemption --- ministerial-approval mandates, by-right upzoning, the builder's remedy --- raises
the marginal cost of a tight \emph{envelope} while leaving the cost of the \emph{discretionary} tax
unchanged. With $r_j = g(\bar z_j, \tau_j)$ and the margins substitutable, a preemption that raises
$\bar z_j$ (the map loosens, as intended) induces a homevoter-pivotal jurisdiction that still demands
restriction to raise $\tau_j$ in compensation: more items conditioned, more opposition entertained,
more continuances. Effective restriction falls by \emph{less} than the change in $\bar z_j$ alone
implies, and --- the sharp testable implication --- \textbf{the conditional approval rate need not
move at all}. Leakage is the share of intended loosening offset on the untouched margin, running from
$0$ (preemption fully effective) to $1$ (fully clawed back). This is the answer to Question 2.

\medskip
\noindent Proposition 3 reframes preemption as a federalism question rather than a mechanical policy:
partial centralization fails not because centralization is wrong but because it is \emph{incomplete},
leaving a decentralized margin through which the spillover reasserts itself. Proposition 2 compounds
it --- as preempted jurisdictions claw back via $\tau$, co-submarket neighbors face a smaller price
effect and adjust in turn.

\subsection{Identification and the mapping to data}

Both headline questions are counterfactual, and only \emph{equilibrium} regulation is observed, never
the objective that generated it; recovering $(\phi_j\lambda_j,\, g(\cdot),\, \text{demand} + w(N))$ is
what licenses the counterfactuals. The reaction-function slope is consistent with rivals --- yardstick
competition, common shocks --- so the note routes identification through the \textbf{preemption ramp}:
the staggered, rule-driven timing of state mandates (builder's-remedy exposure following
Housing-Element non-compliance) is a supply-side shock plausibly orthogonal to a locality's
contemporaneous demand shock, moving $\bar z_j$ for reasons other than local taste.

Each primitive has an address in the data: $\bar z_j$ is read from zoning (largely available from
existing LLM-based measurement); $\tau_j$ is built fresh from planning-commission minutes item by
item; the preemption ramp is dated from HCD Housing-Element compliance records (the panel in
\texttt{output/bay\_area\_recon/hcd\_preemption\_panel/}); prices and quantities come from transaction data.
Proposition 2 becomes a spatial-lag regression of a locality's discretionary tax on its co-submarket
neighbors', instrumented with the preemption variation. Proposition 3 becomes a
difference-in-differences around preemption exposure with the \emph{envelope} and the
\emph{discretionary tax} as \emph{separate} outcomes --- predicting the map loosens, discretion
tightens, and the approval rate stays flat. A short panel straddling the 2017--2023 ramp suffices for
both.

\section{Status and open questions}

\begin{enumerate}
\item \textbf{The model is written; nothing is estimated.} \texttt{toy\_model.tex} is a complete
analytical note with three stated propositions, but the propositions are argued rather than proved
from an explicit solved equilibrium, and no parameter has been taken to data. All numbers in the
project so far are data-feasibility counts, not estimates.

\item \textbf{The multi-jurisdiction extension is the binding gap.} The reaction function of
Proposition 2 requires a panel of jurisdictions, and only SF minutes have been processed. The
progress log carries this explicitly as the open build item: construct the multi-jurisdiction
$\tau_j$ panel and integrate $\bar z_j$ for the strategic-interaction layer. Until then Proposition 2
is untestable and Proposition 1's counterfactual has no estimated primitives.

\item \textbf{The by-right versus ministerial definition is unresolved} --- and it is the model's
hinge. Whether a use is principally permitted (built with certainty, no hearing) as against
conditionally permitted or subject to design review determines what $\bar z_j$ even means, and hence
what the leakage test measures. This is a modeling decision, not a measurement; no further data work
resolves it.

\item \textbf{Identification against the rivals is the acknowledged soft spot.} Separating the
spillover from yardstick competition and common shocks is the question the Guren deck puts to the
advisor directly, along with whether the leakage result (Q2) should stand alone as a first paper with
the structural reaction function (Q1) as the extension.

\item \textbf{The dynamic ratchet is deliberately deferred.} A state-dependent version --- regulation
as path-dependent, with persistent or hysteretic preemption effects propagating across neighbors ---
is noted as a separate extension. It would need a long panel and confronts state dependence versus
unobserved heterogeneity. Neither headline question depends on it, which is why the deep-history data
workstream is scoped out.
\end{enumerate}

\end{document}
